\chapter{What HyperSL is, and what it does well}\label{ch:hslwhat}
\def\tbrk{\hfil\penalty50\hfilneg}% local: a line may break before a provenance tag

\begin{onething}
HyperSL is a 9.03-million-parameter masked autoencoder that reads \emph{one pixel's spectrum} at a time, after every
band has been min-max scaled over the scene it came from; we checked it against its own code and checkpoint and
reproduced its paper's clustering result almost to the decimal, so what the next five chapters find is not a broken
download but a good model built for a different job.
\end{onething}

This chapter opens Part~X. The whole part answers one question the team was asked more than once in September and
October 2026: \emph{there is a published spectral foundation model with open weights; why do we not use it?} The
honest answer has two halves. The first half, this chapter, is what the model is and what it is good at, checked
against its source rather than its abstract. The second half, \cref{ch:hslnorm,ch:hsldetect,ch:hslunmix,ch:hslcloud},
is what happened when it was tried on each of the jobs this product needs done, with every number from a recorded
run. \Cref{ch:hslverdict} puts the two halves side by side.

Paths abbreviate. \file{upstream/} is the authors' code as we hold it, byte-identical to their GitHub commit
\code{e56de7e}. It is kept outside the repository, in the folder
\begin{center}\file{C:/Users/FITT/Documents/hypersl_match_data/hypersl_model/}\end{center}
Experiment records that are no longer in the working tree are cited by the commit that holds them, for example
\file{2d893e6:hypersl-validation/summary.md}, which you read with
\code{git show 2d893e6:research/experiments/hypersl-validation/summary.md}.

\begin{analogy}
A concert pianist who learned every piece from recordings that a studio had normalised so that the loudest note of
each recording was exactly as loud as every other recording's loudest note. She plays beautifully, she can tell a
Chopin nocturne from a Bach fugue with her eyes closed, and if you hum her half a phrase she finishes it. Ask her
instead whether \emph{this} hall is louder than last night's, and she cannot say: the one thing every recording
she learned from had removed was the absolute level. HyperSL is the pianist. The level is reflectance.
\end{analogy}

%======================================================================
\section{Why look at an outside model at all}\label{sec:hslwhat:why}

The product needs three things from spectra, and in autumn 2026 each had a weak spot.
\begin{enumerate}
\item \textbf{Naming a material} by comparing a pixel with a laboratory library (USGS splib07, ECOSTRESS):
\cref{ch:material,ch:namer}. Whole-pixel spectral angle named the right man-made family for 6\% of pixels at 30\%
fill\tbrk\tagM{docs/lld/spectral-family-match.md}.
\item \textbf{Detecting anomalies}: Indradhanu's residual, scored by residual Mahalanobis (\cref{ch:residualmaha}).
\item \textbf{Masking cloud}, so that detections on cloud are not reported (\cref{ch:seghead}).
\end{enumerate}
A \emph{foundation model}\index{foundation model} is a network pretrained on a large, broad, unlabelled data set so
that its internal representation can be reused for many tasks with little or no task-specific training. HyperSL
\citep{kong2025hypersl} advertises exactly that for spectra: pretrained ``on over 300 million spectral instances
worldwide'', able to process ``all forms of hyperspectral data across diverse spectral bands and ranges'', with code
and weights public. If any outside model could cover all three weak spots at once, it was this one. That is why it
was evaluated four separate times between 2026-09-20 and 2026-10-14.

\begin{definition}[Embedding, reconstruction, residual]
An \emph{embedding}\index{embedding} is the short vector a network makes of its input (HyperSL's has 256 numbers).
A \emph{reconstruction} $\hat{\vect{x}}$ is the network's attempt to reproduce its input $\vect{x}$ from partial
information. The \emph{residual} $\vect{r}=\vect{x}-\hat{\vect{x}}$ is what the reconstruction got wrong
(\cref{ch:recon}).
\end{definition}

%======================================================================
\section{What it is, read from the code}\label{sec:hslwhat:arch}

\subsection{One pixel in, one pixel out}

HyperSL's forward pass is the method \code{forward(x, wave, mask_ratio)} of \code{SpectralSharedEncoder} in \file{upstream/model.py}.
Its input $x$ has the shape \shape{(B, 1, L)}: $B$ pixels, one channel, $L$ bands. It has no height and no
width.
\textbf{Each pixel is processed alone; no pixel ever sees another.} Keep that sentence in mind through the whole part:
it is the root of \cref{ch:hsldetect}'s result.

\Cref{tab:hslwhat:trace} traces one PRISMA pixel on the 149-band axis the experiments in this part used (the reason
for 149 bands comes in \cref{sec:hslnorm:conv}).

\begin{table}[htbp]
\centering\small
\begin{tabular}{@{}r>{\raggedright\arraybackslash}p{48mm}l>{\raggedright\arraybackslash}p{50mm}@{}}
\toprule
step & operation & shape & where in \file{model.py} \\
\midrule
0 & min-max scaled spectrum and its wavelengths & \shape{(1, 1, 149)}, \shape{(1, 149)} & input, scaled by the caller \\
1 & three 1-D convolutions over band index, kernels 5, 9, 13, stride 1, \code{padding='same'}; BatchNorm; GELU & 3 $\times$ \shape{(1, 256, 149)} & \code{embedding_layer_1..3} \\
2 & concatenate, then a $1\times1$ convolution $768\to256$; BatchNorm; GELU & \shape{(1, 256, 149)} & \code{embedding_layer_4} \\
3 & transpose to tokens: one token per band & \shape{(1, 149, 256)} & \code{x.transpose(2, 1)} \\
4 & add the wavelength code $\sin/\cos$ of $(\lambda-400)/10$ & \shape{(1, 149, 256)} & \code{pos_tokens} \\
5 & keep a random 5\%: \code{int(149 * 0.05)} $=7$ tokens & \shape{(1, 7, 256)} & \code{random_masking} \\
6 & prepend a learned global token & \shape{(1, 8, 256)} & \code{global_token} \\
7 & 7 transformer blocks, then one more layer that also returns its keys and values & \shape{(1, 8, 256)} & \code{encoder_blocks}, \code{transformer_encoder_layer} \\
8 & embedding $\vect{z}$ = global token + mean of the 7 band tokens & \shape{(1, 1, 256)} & \code{aggregate} \\
9 & cross-attention from $\vect{z}$ back onto the 7 kept tokens & \shape{(1, 7, 256)} & \code{corss_attn} \\
10 & append 142 mask tokens, unshuffle, add the wavelength code again & \shape{(1, 149, 256)} & \code{decoder_forward} \\
11 & 3 transformer blocks; LayerNorm & \shape{(1, 149, 256)} & \code{decoder_blocks} \\
12 & 1-D convolution $256\to1$, kernel 3; sigmoid & \shape{(1, 1, 149)} & \code{decoder_layer} \\
\bottomrule
\end{tabular}
\caption{One pixel through HyperSL, shapes from the code. The reconstruction at step 12 covers all 149 bands, visible
and hidden, and passes through a sigmoid, so every output lies between 0 and 1.}
\label{tab:hslwhat:trace}
\end{table}

Three facts in that table matter later.
\begin{itemize}
\item \textbf{It hides bands, not pixels} (step 5). Indradhanu hides whole pixels and rebuilds them from their
neighbours (\cref{ch:masking}); HyperSL hides 142 of a pixel's own 149 bands and rebuilds them from the other 7 of
the \emph{same} pixel. \Cref{ch:hsldetect} is about what that does to anomalies.
\item \textbf{The tokenizer convolves over band index} (step 1): band $i$'s token mixes bands $i-2\dots i+2$, $i-4\dots
i+4$ and $i-6\dots i+6$ whatever their wavelengths. \Cref{sec:hslnorm:conv} measures what that means on a band list
with gaps.
\item \textbf{The output is squashed into $(0,1)$} (step 12). That only makes sense if every input lies in $[0,1]$,
which is what the per-scene min-max scaling of \cref{sec:hslwhat:data} guarantees.
\end{itemize}

\subsection{The wavelength code}

\begin{definition}[Sinusoidal position code]
A \emph{position code}\index{position code} adds to each token a fixed vector that says where the token sits. HyperSL's
says which wavelength a band token stands for. With $p=(\lambda-400)/10$ and frequencies $\omega_j$, entry $2j$ of
the code is $\sin(p\,\omega_j)$ and entry $2j+1$ is $\cos(p\,\omega_j)$, for $j=0,\dots,127$.
\end{definition}

The code computes the frequencies as
\code{div_term = exp(arange(0, 256, 2) * -(log(100.0) / 256))}, that is
\begin{align}
\omega_j &= \exp\!\left(2j\cdot\left(-\frac{\ln 100}{256}\right)\right) && \text{the line above, with }2j\text{ for }\texttt{arange(0,256,2)}\notag\\
&= \exp\!\left(\ln 100\right)^{-2j/256} && \text{since }e^{-ab}=(e^{a})^{-b}\notag\\
&= 100^{-2j/256}. \label{eq:hslwhat:omega}
\end{align}
So $\omega_0=100^0=1$, $\omega_1=100^{-2/256}=0.964662$, and $\omega_{127}=100^{-254/256}=0.010366$\tagC.

\begin{worked}[label={we:hslwhat:pe}]{The wavelength code at 450\,nm}
\begin{align*}
p &= (450-400)/10 = 50/10 = 5.0,\\
\text{entry }0 &= \sin(5.0\times1) = \sin 5.0 = -0.9589,\\
\text{entry }1 &= \cos(5.0\times1) = \cos 5.0 = +0.2837,\\
\text{entry }254 &= \sin(5.0\times0.010366) = \sin 0.05183 = +0.0518 .
\end{align*}
The fastest pair turns through $5.0$ radians, the slowest through $0.052$: the slow entries tell 450\,nm from 2400\,nm
($p=200$, entry 254 is $\sin 2.0732=+0.8764$), the fast ones tell neighbours apart\tagC. Any wavelength has a code,
so the model can be fed any band list.
\end{worked}

\begin{sourcenote}[: the base is 100 in the code, 10,000 in the paper]
The paper's equations (5) and (6) divide $w_i-400$ by $10\,000^{2j/d}$ \citep{kong2025hypersl}. The code divides by
\code{100} raised to that power (\code{math.log(100.0)}). The released weights were trained with the code, so the
code is right for them; feeding the paper's frequencies to the checkpoint would be off-distribution. A second
check: the wavelengths must be in nanometres. In micrometres, 1.0\,$\mu$m gives $p=(1.0-400)/10=-39.9$, and the
validation run measured the loss rising from 0.0024 to 0.173 on EnMAP\tbrk\tagMr{\texttt{2d893e6}: summary.md, step 2,
row G}.
\end{sourcenote}

\begin{sourcenote}[: the last encoder layer has no skip connection]
\code{TransformerEncoderLayer.forward} reads \code{x, k, v = self.attn(self.norm1(x))} and then
\code{x = x + self.drop_path(x)}. The second \code{x} is the attention output, not the layer's input, so the line
doubles the attention output and the usual residual (skip) connection of \cref{ch:block} is lost in this one layer.
The weights were trained this way, so ``fixing'' it would break the checkpoint. It does not change any conclusion in
this part; it is the kind of thing you only find by reading the source.
\end{sourcenote}

\subsection{Counting the parameters, by hand}

\begin{worked}[label={we:hslwhat:params}]{HyperSL's 9,027,841 trainable parameters}
A linear layer from $i$ to $o$ numbers has $i\times o$ weights and $o$ biases; a LayerNorm or BatchNorm over $D$
numbers has $2D$ (a scale and a shift). With $D=256$ and the feed-forward width $4D=1024$, one transformer block is
\begin{align*}
\text{two LayerNorms} &= 2\times(2\times256) = 1{,}024,\\
\text{query, key, value} &= 256\times768+768 = 196{,}608+768 = 197{,}376,\\
\text{output projection} &= 256\times256+256 = 65{,}792,\\
\text{feed-forward} &= (256\times1024+1024)+(1024\times256+256)\\
 &= 263{,}168+262{,}400 = 525{,}568,\\
\text{block} &= 1{,}024+197{,}376+65{,}792+525{,}568 = 789{,}760 .
\end{align*}
The encoder has 8 such layers and the decoder 3:
\begin{align*}
\text{encoder} &= 8\times789{,}760 = 6{,}318{,}080,\\
\text{decoder} &= 3\times789{,}760 = 2{,}369{,}280,\\
\text{cross-attention} &= 2\times(256\times256+256) = 131{,}584,\\
\text{tokenizer convolutions} &= \underbrace{(5\cdot256+256)}_{1{,}536}+\underbrace{(9\cdot256+256)}_{2{,}560}
 +\underbrace{(13\cdot256+256)}_{3{,}584}\\
 &\quad+\underbrace{(768\cdot256+256)}_{196{,}864} = 204{,}544,\\
\text{tokenizer BatchNorms} &= 4\times(2\times256) = 2{,}048,\\
\text{global and mask tokens} &= 2\times256 = 512,\\
\text{two final LayerNorms} &= 2\times(2\times256) = 1{,}024,\\
\text{output convolution} &= 256\times3+1 = 769,\\
\text{total} &= 6{,}318{,}080+2{,}369{,}280+131{,}584+204{,}544\\
&\quad+2{,}048+512+1{,}024+769\\
&= 9{,}027{,}841 .
\end{align*}
The released checkpoint holds 9,158,021 numbers. The difference, $9{,}158{,}021-9{,}027{,}841=130{,}180$, is exactly
its stored buffers: the position table \code{pe} ($500\times256=128{,}000$), \code{div_term} (128), and the
BatchNorms' running means and variances ($4\times2\times256=2{,}048$) with their four step counters
($128{,}000+128+2{,}048+4=130{,}180$)\tagC. The ``$\sim$9.16\,M parameters'' quoted in the team's notes counts the
buffers too.
\end{worked}

The detection run of \cref{ch:hsldetect} logged the model's \code{parameters()} as 9,155,969\tbrk\tagM{research/experiments/hypersl-anomaly-bench/results.json}: the trainable 9,027,841 plus
\code{pe} and \code{div_term}, which the code declares as frozen parameters rather than buffers ($9{,}027{,}841+128{,}000+
128=9{,}155{,}969$)\tagC. The tokenizer is $(204{,}544+2{,}048)/9{,}027{,}841=2.29\%$ of the model, and 96\% of the tokenizer is its $1\times1$
fusion convolution: the part that ``reads the band list'' is tiny\tagC.

\begin{sourcenote}[: the class defaults are not the checkpoint]
\code{SpectralSharedEncoder}'s defaults are \code{embedding_dim=1024, encoder_depth=32, decoder_depth=16}. The
checkpoint loads only into 256 / 8 / 4 with 8 heads, the configuration \file{upstream/train.py} builds. A strict
state-dict load checks the sizes but not the head count, which was confirmed by behaviour: 4 or 16 heads load
without error and raise the EnMAP test loss from 0.0024 to 0.087 or 0.008\tbrk\tagMr{\texttt{2d893e6}: summary.md,
step 2}. Every upstream downstream script ships with its checkpoint-load lines commented out, so running one
unmodified evaluates a randomly initialised network.
\end{sourcenote}

%======================================================================
\section{What it was trained on}\label{sec:hslwhat:data}

\subsection{The data path, as code}

\file{upstream/mutihsi.py} builds the pretraining file. For an EnMAP scene, \code{read_EnMap} reads the L2A
reflectance, multiplies by $10^{-4}$, deletes band indices 128 to 134 (the paper's ``bad bands [128--134]''), sets
nodata to NaN, reshapes the cube to a table of pixels by bands, and calls \code{self.scale}, which is
\code{sklearn.preprocessing.minmax_scale}. DESIS loses bands 0 to 5. The scaling is applied per \emph{column} of that
table, which is per band, over every pixel of the scene:
\begin{equation}
x'_{b}=\frac{x_b-\min_b}{\max_b-\min_b},\qquad \min_b,\ \max_b \text{ taken over the scene's pixels in band } b.
\label{eq:hslwhat:minmax}
\end{equation}
\Cref{ch:hslnorm} takes this formula apart. Here note only what it means for training: \textbf{HyperSL has never seen
a reflectance.} Every spectrum it learned from was rescaled so that, in each band, the darkest pixel of its own scene
was 0 and the brightest 1. Water-absorption bands were kept: 217 of EnMAP's 224 bands went in, and all but the first 6 of DESIS's.

\begin{sourcenote}[: the paper describes the scaling in four words]
The paper says ``to ensure that the features of each band are on the same scale, we performed layer-by-layer
normalization'' \citep{kong2025hypersl}. It does not say min-max, and it does not say per scene; both are visible only
in the code (\file{mutihsi.py}, and again in every \file{downstream_*.py} reader). The team's earlier note that ``the
paper does not mention it'' is slightly too strong: it mentions a normalisation, not which one.
\end{sourcenote}

\subsection{What the training set is, and what it is not}

\begin{center}\small
\begin{tabular}{@{}>{\raggedright\arraybackslash}p{30mm}>{\raggedright\arraybackslash}p{106mm}@{}}
\toprule
claim & what the code and files show \\
\midrule
sensors & EnMAP and DESIS L2A only (\file{train.py} builds two loaders, \code{'ENMap'} and \code{'DESIS'}). A GF5 reader
exists in \file{mutihsi.py} and is never called. The paper's 439-band Shanghai airborne set is not in \file{train.py}. \\
size & ``over 300 million spectral instances'' (abstract); ``approximately 200 million spectral data points'' (section
IV-A). The team's reading: about 300 one-megapixel tiles\tbrk\tagMr{memory note on HyperSL, 2026-10-04}. \\
clouds & ``cloud cover near zero''; ``overall cloud coverage of the data is below 10\%'' (paper, Fig.~6 caption).
\Cref{ch:hslcloud} needs this. \\
training schedule & 10 epochs, AdamW $5\times10^{-4}$, batch 256 per sensor, mask ratio 0.95, a checkpoint every 5
epochs (\file{train.py}); the released \code{10_base_mask95} is epoch 10\tbrk\tagC. \\
loss & $\text{MSE}+(1-\cos)$ over the \emph{whole} reconstructed spectrum, visible bands included
(\file{loss.py}, \code{MSE_SAM_loss}). \\
where the data is & the $\sim$265\,GB training corpus is on Baidu Netdisk only; Hugging Face has the 10.5\,GB
\emph{test} split (about 12 million spectra)\tbrk\tagMr{memory note on HyperSL}. \\
licence & no LICENSE file in the GitHub repository; the Hugging Face card says MIT\tbrk\tagMr{memory note on
HyperSL, ``verify before shipping''}. \\
\bottomrule
\end{tabular}
\end{center}

%======================================================================
\section{Did we have the right model? The reproduction}\label{sec:hslwhat:repro}

Before blaming a model you must show you ran it right. The validation (2026-10-02/03) did that in four steps; all
numbers below are from \texttt{2d893e6:hypersl-validation/summary.md}.

\subsection{Step 1: provenance}

The checkpoint's sha256 is \code{3d5f1914...a74a68c}, identical to Hugging Face's \code{10_base_mask95_checkpoint.pt};
the model, data and downstream scripts are byte-identical to upstream \code{e56de7e} (2026-02-03); the configuration
256 / 8 / 4 / 8 heads matches \file{train.py}. This book's own script re-hashed the file and read its
\code{epoch} field: 10\tbrk\tagC.

\subsection{Step 2: the right settings, by elimination}

There is no published loss on the test split, so the validation ran nine variants on 20,000 test spectra per sensor
(loss $\text{MSE}+(1-\cos)$ at mask 0.95, mean of 3 mask seeds; lower is better):

\begin{center}\small
\begin{tabular}{@{}llllll@{}}
\toprule
& variant & DESIS loss & EnMAP loss & DESIS SAM & EnMAP SAM \\
\midrule
A & \textbf{ours: epoch 10, 8 heads, nm} & \textbf{0.0050} & \textbf{0.0024} & \textbf{4.4°} & \textbf{2.8°} \\
B & epoch 5 & 0.0054 & 0.0032 & 4.7° & 3.3° \\
C & epoch 0 & 0.0066 & 0.0053 & 5.5° & 4.1° \\
D & untrained & 0.183 & 0.218 & 25.4° & 26.7° \\
E & 4 heads & 0.054 & 0.087 & 16.2° & 19.2° \\
F & 16 heads & 0.012 & 0.008 & 7.4° & 5.7° \\
G & wavelengths in $\mu$m & 0.132 & 0.173 & 24.1° & 27.9° \\
H & wavelengths $+20$\,nm & 0.0081 & 0.0073 & 5.9° & 5.3° \\
I & per-spectrum 0--1 scaling & 0.0137 & 0.0063 & 6.8° & 4.1° \\
\bottomrule
\end{tabular}
\end{center}
Row A wins every column, and epochs 0, 5, 10 improve in order: the setup is right. Row I matters for
\cref{ch:hslnorm}: rescaling each spectrum by its \emph{own} minimum and maximum instead of the scene's roughly doubles
the loss. A second finding from the same run: with mask ratio 0 (all bands visible) the loss is \emph{worse} than at
0.95 (0.0081 / 0.0073 against 0.0050 / 0.0024), because the model only ever trained at 0.95.

\subsection{Step 3: the paper's own benchmarks}

\paragraph{Zero-shot reconstruction (paper Table III).} On Indian Pines, an AVIRIS scene HyperSL never saw, the
reproduction was equal to or better than the paper on every measure (SAM 7.86° at mask 0.95 against the paper's
8.43°). On Botswana (Hyperion) MSE was ten times better but SAM 2--3° worse (8.84° against 5.67°), most likely
because the authors' wavelength file for Botswana is unpublished.

\paragraph{Unsupervised classification on Chikusei (paper Table XII).} Chikusei is a labelled 19-class airborne scene.
Embed every labelled pixel (77,592), cluster the embeddings with K-means into 19 groups, map clusters to classes
one-to-one by the best assignment, and score:

\begin{definition}[Overall accuracy, adjusted Rand index]
\emph{Overall accuracy}\index{overall accuracy} (OA) is the share of pixels whose mapped cluster equals their class.
The \emph{adjusted Rand index}\index{adjusted Rand index} (ARI) compares two groupings pair by pair, corrected so that a
random grouping scores 0 and an identical one 1 \citep{hubert1985comparing}.
\end{definition}

\begin{center}\small
\begin{tabular}{@{}lll@{}}
\toprule
features + K-means & OA (\%) ours / paper & ARI ours / paper \\
\midrule
raw spectra & 39.60 / 42.69 & 0.268 / 0.289 \\
PCA & 37.68 / 41.11 & 0.259 / 0.276 \\
\textbf{HyperSL (band 0--1 input)} & \textbf{59.60 / 60.23} & \textbf{0.5404 / 0.5405} \\
HyperSL (raw input) & 48.19 / -- & 0.418 / -- \\
\bottomrule
\end{tabular}
\end{center}
HyperSL's row matches the paper to 0.63 OA points and 0.0001 ARI\tagC. Look at the last row: the same embeddings from
\emph{raw} reflectance lose $59.60-48.19=11.41$ points\tagC. The authors' commented-out recipe in their script feeds
raw values, so ``the paper's numbers were not produced by that code as written''; they need the per-band scaling.
Two rows did not reproduce: agglomerative clustering on HyperSL features reached 62.23 against the paper's 73.43,
unexplained; spectral clustering is not comparable because its affinity is unpublished.

\begin{figure}[htbp]
\centering
\begin{tikzpicture}
\begin{axis}[ybar, bar width=9pt, width=0.85\linewidth, height=55mm, ymin=0, ymax=70,
  symbolic x coords={raw,PCA,HyperSL 0-1,HyperSL raw}, xtick=data, ylabel={OA (\%)},
  legend style={at={(0.02,0.98)},anchor=north west,font=\small}, nodes near coords, nodes near coords style={font=\tiny},
  enlarge x limits=0.18]
\addplot[fill=spat!60] coordinates {(raw,39.60) (PCA,37.68) (HyperSL 0-1,59.60) (HyperSL raw,48.19)};
\addplot[fill=maskgrey!70] coordinates {(raw,42.69) (PCA,41.11) (HyperSL 0-1,60.23)};
\legend{ours (2026-10-02), paper Table XII}
\end{axis}
\end{tikzpicture}
\caption[Chikusei clustering: HyperSL features against raw spectra and PCA]{Chikusei, K-means on 19 clusters.
HyperSL's features add 20 points of overall accuracy over the raw spectra, as the paper says, but only when the input is
scaled per band; on raw reflectance the gain halves. Numbers from
\texttt{2d893e6:\allowbreak hypersl-validation/\allowbreak summary.md}.}
\label{fig:hslwhat:chikusei}
\end{figure}

\paragraph{Fine-tuning (paper Tables VI and VII).} Ten seeds on Colab, the paper's protocol (patch 5, 10\% of each class,
150 epochs, AdamW $10^{-4}$, weight decay 0.05):

\begin{center}\small
\begin{tabular}{@{}lll@{}}
\toprule
data, initialisation & OA ours & OA paper \\
\midrule
Botswana, pretrained & $99.01\pm0.51$ & $99.66\pm0.22$ \\
Botswana, random & $97.55\pm2.13$ & $97.05\pm1.48$ \\
Indian Pines, pretrained & $97.73\pm0.52$ & $98.48\pm0.39$ \\
Indian Pines, random & $94.94\pm3.05$ & $95.04\pm1.78$ \\
\bottomrule
\end{tabular}
\end{center}

\begin{worked}[label={we:hslwhat:gain}]{How much does pretraining buy?}
The gain is pretrained minus random.
\begin{align*}
\text{Botswana, ours} &= 99.01-97.55 = +1.46, & \text{paper} &= 99.66-97.05 = +2.61,\\
\text{ratio} &= 1.46/2.61 = 0.56,\\
\text{Indian Pines, ours} &= 97.73-94.94 = +2.79, & \text{paper} &= 98.48-95.04 = +3.44,\\
\text{ratio} &= 2.79/3.44 = 0.81 .
\end{align*}
Pretraining helps on both, as claimed, by 56\% and 81\% of the paper's gain\tagC. The random-initialisation rows match
the paper, so data, split and head agree with the authors'; the remaining gap sits in unpublished fine-tuning details.
(The summary's own rounding, ``+1.45'' and ``+2.80'', differs from these by 0.01 in the last digit.)
\end{worked}

%======================================================================
\section{What it is good at}\label{sec:hslwhat:good}

Put the reproduction together and HyperSL's strengths are real and specific.
\begin{itemize}
\item \textbf{It learns the shape of spectra.} On a scene it never saw, from a sensor it never saw (Indian Pines), it
reconstructs 95\% hidden bands at 7.9° from seven visible ones.
\item \textbf{Its embedding groups land cover without labels}, 20 points better than raw spectra on Chikusei.
\item \textbf{It is a good starting point for supervised classification}, worth 1.5 to 2.8 OA points over random
initialisation on two classic scenes.
\item \textbf{It accepts any band list}, because its position code is a function of wavelength, not of band number.
Our own stem has the same property (\cref{ch:crossstem,ch:specdecoder}); the idea is sound.
\end{itemize}
All four strengths share a setting: \emph{one scene at a time, classes that cover many pixels, and a frame in which only
relative differences inside that scene matter}. Classification of a labelled benchmark scene is exactly that setting.

\begin{keypoint}
HyperSL is well built for describing land cover inside one scene. Each of our three jobs asks something it was not
built for: absolute reflectance that can be compared with a laboratory (library matching), surprise at a few pixels
that do not fit their surroundings (anomaly detection), and recognising a class its training data excluded (cloud).
\end{keypoint}

\begin{center}\small
\begin{tabular}{@{}>{\raggedright\arraybackslash}p{30mm}>{\raggedright\arraybackslash}p{52mm}>{\raggedright\arraybackslash}p{52mm}@{}}
\toprule
our job & what it needs from a model & HyperSL's design \\
\midrule
library matching & the same spectrum gives the same input in every scene, and so does a library spectrum &
input depends on the scene's extremes; a library spectrum has no scene (\cref{ch:hslnorm}) \\
anomaly detection & a pixel must not help predict itself; context from neighbours &
predicts a pixel's bands from its own bands; no neighbours (\cref{ch:hsldetect}) \\
unmixing, sub-pixel naming & per-pixel abundances that keep a 0.1-pixel target &
deep unmixing head smooths over $3\times3$ (\cref{ch:hslunmix}) \\
cloud masking & clouds in the training data, brightness kept &
trained on near-cloud-free scenes, brightness rescaled away (\cref{ch:hslcloud}) \\
\bottomrule
\end{tabular}
\end{center}

\begin{drift}
The checkpoint (\code{10_base_mask95}, epoch 10), the upstream commit \code{e56de7e}, the licence status, the test split
on Hugging Face and the reproduction numbers are as of 2026-10-04. A later HyperSL release could change any of them.
\end{drift}

%======================================================================
\begin{confusions}
\confusion{HyperSL looks at image patches like Indradhanu.}{It looks at one pixel's spectrum. Its ``tokens'' are bands,
and it never sees a neighbour. (Its classification fine-tune uses $5\times5$ patches, but that is a downstream head, not
the pretrained model.)}
\confusion{It has 9.16 million parameters.}{9,027,841 trainable; 9,158,021 counting stored buffers, of which 128,000 is
the position table.}
\confusion{The wavelength code makes the tokenizer wavelength-aware.}{The code is added \emph{after} the convolutions. The
convolutions themselves mix neighbouring band indices, whatever their wavelengths (\cref{sec:hslnorm:conv}).}
\confusion{Mask 0 (show every band) gives the best embedding.}{The model never trained at mask 0; its loss there is worse
than at 0.95.}
\confusion{The reproduction proves HyperSL suits our product.}{It proves we ran the model the authors trained, as they
ran it. Whether it suits a different job is a separate question, answered in the next four chapters.}
\confusion{The paper hides the normalisation.}{It names ``layer-by-layer normalization''; the code shows it is min-max per
band per scene.}
\end{confusions}

\begin{checkyourself}
\item\label{ex:hslwhat:1} A HyperSL with embedding width $D=128$ and the same depths (8 encoder layers, 3 decoder blocks, feed-forward
  $4D$). Count one transformer block's parameters.
\item\label{ex:hslwhat:2} Compute entries 0 and 1 of the wavelength code at 1000\,nm, and say what $p$ would be if someone passed the
  wavelength as 1.0 (micrometres).
\item\label{ex:hslwhat:3} EnMAP L2A has 224 bands and HyperSL drops 7. At mask ratio 0.95, how many bands does the encoder see? How many
  for SpectralEarth's 202-band list?
\item\label{ex:hslwhat:4} On Chikusei, by how many points does HyperSL beat raw spectra with band 0--1 input, and with raw input? What
  share of the first gain survives without the scaling?
\item\label{ex:hslwhat:5} The checkpoint holds 9,158,021 numbers. Without looking, list the four kinds of stored buffer and their sizes,
  and check they sum to the difference from 9,027,841.
\item\label{ex:hslwhat:6} Why does a sigmoid at the output require the per-scene min-max at the input? What would happen to a reflectance
  of 1.2 (a specular glint) fed raw?
\end{checkyourself}

\nextchapter{ch:hslnorm}
